\subsection{Group homomorphisms}
In this section, we will define the concept of group homomorphism, which we 
will later use to proof a result about the solvability of polynomial equations.  

\begin{definition}[Group homomorphism]
    A group homomorphism is a function $\varphi: G \to H$ between 
    two groups $G$ and $H$ such that:
    
    $$ \forall g_1, g_2 \in G\;\; \varphi(g_1g_2) = \varphi(g_1)\varphi(g_2) $$
\end{definition} \vspace{0.2em}

\begin{definition}[Kernel]
    We call kernel of a group homomorphism $\varphi: G \to H$ the set of
    elements in $G$ that are mapped to the identity element of $H$ by $\varphi$.

    $$ \texttt{ker}_{\varphi} = \{ g \in G \;:\; \varphi(g) = 1_H \}$$
\end{definition} \vspace{0.2em}

\begin{definition}[Group epimorphism]
    A group homomorphism $\varphi: G \to H$ that is also a surjective 
    function is called a group epimorphism.
\end{definition} \vspace{0.2em}

\begin{lemma}
\label{lem:epimorphismPreservesNormality}
    The image of $H$ normal subgroup of $G$ via the epimorphism $\varphi: G \to K$
    is a normal subgroup of the image of $K$ via the same epimorphism 
    ($H \triangleleft G \Rightarrow \vv{\varphi}(H) \triangleleft K$). 
\end{lemma}
\begin{proof}
    Consider the group $\vv{\varphi}(H)$ and show that it is normal in $K = \vv{\varphi}(G)$:

    $$ \forall k \in K\;\; \forall l \in \vv{\varphi}(H)\;\; klk^{-1} = 
        \varphi(g)\varphi(h)\varphi(g^{-1}) = \varphi(ghg^{-1}) \in \vv{\varphi}(H) $$ \\

    This proof uses the fact that $\varphi$ is an epimorphism, hence it is surjective, to 
    conclude that for every $k,l$ exists $g, h$ such that $klk^{-1} = \varphi(g)\varphi(h)
    \varphi(g^{-1})$. For non suriective homomorphisms, the proof is not valid, because 
    there might be elements in $K$ that are not images of elements in $G$, herby 
    it could be impossible to find $g,h$. \\
\end{proof} \vspace{0.2em}

\begin{definition}[Group endomorphism]
    A group homomorphism $\varphi: G \to H$ that is also a injective 
    function is called a group endomorphism.
\end{definition} \vspace{0.2em}

\begin{definition}[Group isomorphism]
    A group homomorphism $\varphi: G \to H$ that is both injective and surjective 
    (herby a bijection) is called a group isomorphism.
\end{definition} \vspace{0.2em}

\begin{lemma}
\label{lem:abelianityPreservedByIsomorphism}
    If $A$ is an abelian group and $\varphi: A \to B$ is a group isomorphism, 
    then $B$ is also an abelian group (notice that since a bijection is surjiective 
    $B = \vv{\varphi}(A)$).
\end{lemma}
\begin{proof}
    Since $\varphi$ is a group isomorphism, it is a bijection, herby every 
    element in $b_i \in B$ is the image of an element in $a_i \in A$ via $\varphi$, 
    thus, the following holds true:

    $$ b_1b_2 = \varphi(a_1)\varphi(a_2) = \varphi(a_1a_2) = \varphi(a_2a_1) = 
        \varphi(a_2)\varphi(a_1) = b_2b_1 $$ \\

\end{proof} \vspace{0.2em}

\newpage

\begin{theorem}[First theorem of isomorphism]
\label{thm:firstIsomorphism}
    For every group \emph{epimorphism} $\varphi: A \to B$, the image of $A$ via $\varphi$ 
    is isomorphic to the quotient group $A$ divided by the kernel of $\varphi$:

    $$ \frac{G}{\texttt{ker}_{\varphi}} \sim \vv{\varphi}(G) $$ \vspace{0.2em}

    We can proof that these groups are in fact isomorphic by showing that the
    function $\psi$ is an isomorphism between them, since it's clearly both 
    injective and surjective.

    $$ \psi: \frac{G}{\texttt{ker}_{\varphi}} \to \vv{\varphi}(G) $$
    $$ \psi(g\texttt{ker}_{\varphi}) = \varphi(g) $$ \vspace{0.2em}
\end{theorem}

\begin{theorem}[Second theorem of isomorphism]
\label{thm:secondIsomorphism}
    Consider the groups $H \triangleleft B \triangleleft A$. Then, there is an 
    isomorphism between the quotient groups $(AH/BH)$ and $\frac{A}{B}$.

    $$ \frac{AH}{BH} \sim \frac{A}{B} $$ \vspace{0.2em}
\end{theorem}
\begin{proof}
    To prove that exists an isomorphism between such groups, we will rely on the previous 
    result of the first theorem of isomorphism (theorem \ref{thm:firstIsomorphism}). We 
    therefore need to exibit a homomorphism between the groups $AH$ and $\frac{A}{B}$:


    $$ \varphi : (AH) \rightarrow \frac{A}{B} $$
    $$ \varphi(ah) = aB $$ \vspace{0.2em}

    Constructed as such, we can see that $\varphi$ is a homomorphism, since $B$ is a normal subgroup 
    of $A$, hence, left cosets and right cosets are equal (lemma \ref{lem:normal_cosets}).

    $$ \varphi(a_1h_1)\varphi(a_2h_2) = a_1Ba_2B = a_1B(a_2B) = a_1a_2BB = a_1a_2B = \varphi(a_1h1a_2h_2) $$ \\

    We can also see that it is a surjective function, since for every $aB$ in $\frac{A}{B}$,
    we can find $ah$ in $AH$ such that $\varphi(ah) = aB$. So let's go ahead and inspect 
    it's kernel (the set of elements mapped to the identity). Please note that the kernel 
    of this epimorphism is just $BH$: \\

    $$ \varphi(ah)\varphi(b_1h_1 \in BH \subseteq AH) = aB = \varphi(ah) $$
    $$ \texttt{ker}_{\varphi} = BH $$

    We can herby apply the first theorem of isomorphism (theorem \ref{thm:firstIsomorphism}) to
    finish the proof:

    $$ \frac{AH}{BH} = \frac{AH}{\texttt{ker}_{\varphi}} \sim \vv{\varphi}(AH) = \frac{A}{B} $$ \\

\end{proof} \vspace{0.2em}

\newpage

\begin{theorem}[Third theorem of isomorphism]
\label{thm:thirdIsomorphism}
    Consider the groups $H \triangleleft B \triangleleft A$. Then, there is an 
    isomorphism between the quotient groups $(\frac{A}{H}/\frac{B}{H})$ and $\frac{A}{B}$.

    $$ \frac{(\frac{A}{H})}{(\frac{B}{H})} \sim \frac{A}{B} $$ \vspace{0.2em}
\end{theorem}
\vspace{0.2em}

To prove this theorem, it's helpful to keep in mind the form of the generic element of 
the set (groups) we're working with (for an easier notation, we will call $H\frac{B}{H}$ by 
the name $K$):

$$ \frac{A}{H} = \{ aH \;:\; a \in A \} $$
$$ \frac{B}{H} = \{ bH \;:\; b \in B \} $$
$$ \frac{A}{B} = \{ aB \;:\; a \in A \} $$
$$ \frac{(\frac{A}{H})}{(\frac{B}{H})} = \{ aK \;:\; a \in A \} $$ 
\vspace{0.2em}

\begin{proof}
    To prove that exists an isomorphism between such groups, we will just exibit it:

    $$ \varphi : (\frac{A}{H})/(\frac{B}{H}) \rightarrow \frac{A}{B} $$
    $$ \varphi(aK) = aB $$ \vspace{0.2em}

    First things first: let's show that it is indeed an homomorphism (compatibility with $\cdot$):

    $$ 
        \varphi(rK)\varphi(tK) = 
        rBtB = r(Bt)B = r(tB)B = rt(BB) = rtB = (rt)B = 
        \varphi(rtK)
    $$ \vspace{0.2em}

    Note that here in this very argument we used the all these subgroups are normal, hence 
    every right coset is equal to a left coset (lemma \ref{lem:normal_cosets}) when writing 
    $ r(Bt)B = r(tB)B $. \\


    We can clearly see that it is a bijection, since for every $aB$ in $\frac{A}{B}$, 
    we can find $aK$ in $(\frac{A}{H})/(\frac{B}{H})$ such that $\varphi(aK) = aB$. Thus 
    we only really need to show that it's an injection to complete the proof. Injectivity can 
    be shown with the following argument: \\

    $$ 
        \varphi(rK) = \varphi(tK) \Rightarrow rB = tB \Rightarrow r\epsilon = 
        t\epsilon \Rightarrow r = t 
    $$ \vspace{0.2em}

    These key obeservations show that the function $\varphi$ is indeed a bijective 
    homomorphism, herby it is an isomorphism. To express the fact that an isomorphism exists
    betweek two groups the symbol $\sim$ is used 
    (sometimes it can be used $\cong$ as well, but not here). \\ 

\end{proof} \vspace{0.2em}

\newpage